\documentclass[11pt,a4paper]{article}
\usepackage[margin=1in]{geometry}
\usepackage{amsmath,amssymb,amsthm}
\usepackage{hyperref}

\newtheorem{theorem}{Theorem}
\newtheorem{definition}{Definition}
\newtheorem{remark}{Remark}

\title{N=3 Stationary Geometry and Geometric Phase\\from the Connectivity Operator \(C(r)\)}
\author{Derived from the Two-Rocks / EGT Operator}
\date{}

\begin{document}
\maketitle

\begin{abstract}
We consider three distinguishable elements (``rocks'') whose pairwise interactions are governed by the fixed connectivity operator
\[
C(r)=0.001\,(1+2r)\,e^{-r/3}\,e^{i(\pi/4)r}.
\]
When the relational distances are allowed to evolve by a flow that increases with \(|C(r)|\), the unique stationary configuration is the equilateral geometry with all sides equal to the preferred scale \(r=2.5\). Around the closed triangular loop the accumulated geometric phase evaluates to \(15\pi/8\). Local attractiveness of the stationary point follows from the negative second derivative of \(\ln|C(r)|\) at \(r=2.5\). No additional postulates beyond the operator, conservation of total amplitude, and the dynamical-distance rule are used.
\end{abstract}

\section{Primitives}
\begin{definition}[Connectivity operator]
\[
C(r)=0.001\,(1+2r)\,e^{-r/3}\,e^{i(\pi/4)r},\qquad r\ge0.
\]
\end{definition}

The magnitude \(|C(r)|\) possesses a unique maximum at
\[
r_{\mathrm{opt}}=2.5,
\]
obtained by solving
\[
\frac{d}{dr}\ln|C(r)|=\frac{2}{1+2r}-\frac13=0.
\]

\begin{definition}[N=3 state and generator]
Let the three rocks be labeled \(\{1,2,3\}\) with relational distances \(r_{12}=x\), \(r_{13}=y\), \(r_{23}=z>0\). A state is a vector \(\psi\in\mathbb{C}^3\). The generator \(A\) is the \(3\times3\) matrix
\begin{align*}
A_{ij}&=C(r_{ij})\qquad(i\neq j),\\
A_{ii}&=-\sum_{k\neq i}C(r_{ik}).
\end{align*}
The diagonal implements conservation of total amplitude \(\sum_k\psi_k\).
\end{definition}

\section{Dynamical distances}
We allow the distances themselves to evolve. The minimal rule consistent with the operator is a flow that increases the pairwise weight when the distance is away from the maximum of \(|C|\):
\[
\dot r_{ij}\;\propto\;-\frac{\partial}{\partial r_{ij}}\bigl(|C(r_{ij})|^2\bigr).
\]
Any flow monotonic in \(|C|\) drives each pairwise distance toward \(r=2.5\).

\section{Stationary geometry}
\begin{theorem}[Equilateral stationary point]
Under the flow above, the unique stationary configuration is
\[
x=y=z=2.5.
\]
\end{theorem}

\begin{proof}
Stationarity requires \(\dot x=\dot y=\dot z=0\). This occurs if and only if each distance sits at the unique critical point of \(|C|\), i.e., at \(r=2.5\). Thus the only stationary geometry is the equilateral triangle of side length \(2.5\).
\end{proof}

\section{Geometric phase}
\begin{theorem}[Loop phase]
Around the closed ordered loop \(1\to2\to3\to1\) the total geometric phase is
\[
\frac{\pi}{4}(x+y+z).
\]
At the stationary point this evaluates to
\[
\frac{\pi}{4}\cdot7.5=\frac{15\pi}{8}.
\]
\end{theorem}

\begin{proof}
Each transfer across a link of length \(r\) multiplies the amplitude by the phase factor \(e^{i(\pi/4)r}\). The product around the cycle is therefore \(e^{i(\pi/4)(x+y+z)}\). Substituting the stationary values yields the stated result.
\end{proof}

\section{Local stability}
\begin{theorem}[Attractivity]
The equilateral point \(x=y=z=2.5\) is locally attractive under the gradient flow of \(|C|^2\).
\end{theorem}

\begin{proof}
The second derivative of \(\ln|C(r)|\) is strictly negative at \(r=2.5\). Consequently each pairwise coordinate has a restoring force toward \(2.5\). The constrained three-dimensional flow on \((x,y,z)\) therefore admits the equilateral point as a local attractor.
\end{proof}

\section{Remarks}
\begin{remark}
All results follow from the single operator \(C(r)\), conservation of total amplitude, and the decision to let distances evolve by a flow monotonic in \(|C|\). No external physical constants or additional postulates are introduced.
\end{remark}

\begin{remark}
The absolute conversion from the dimensionless scale \(r=2.5\) into a laboratory length (or frequency) is not fixed by the operator and remains an identification to be supplied by experiment.
\end{remark}

\begin{remark}
The present note derives only the stationary geometry and the loop phase. Any claim that these quantities constitute a measured entanglement witness, a double-slit visibility, or any other laboratory observable requires further explicit mapping steps that lie outside this document.
\end{remark}

\end{document}
